% Generated draft: requires mathematical and source review.
\documentclass[11pt]{article}
\usepackage[margin=25mm]{geometry}
\usepackage{fontspec}
\setmainfont{Latin Modern Roman}
\usepackage{amsmath,amssymb,mathtools,unicode-math}
\setmathfont{Latin Modern Math}
\usepackage{enumitem,needspace,xurl,hyperref}
\setlength{\parindent}{0pt}
\setlength{\parskip}{5pt}
\setlength{\emergencystretch}{3em}
\newcommand{\N}{\mathbb{N}}
\newcommand{\Z}{\mathbb{Z}}
\newcommand{\Q}{\mathbb{Q}}
\newcommand{\R}{\mathbb{R}}
\newcommand{\C}{\mathbb{C}}
\newcommand{\F}{\mathbb{F}}
\newcommand{\A}{\mathbb{A}}
\newcommand{\PP}{\mathbb{P}}
\newcommand{\E}{\mathbb{E}}
\newcommand{\eps}{\varepsilon}
\newcommand{\dd}{\,\mathrm{d}}
\begin{document}
\section*{UnsolvedMath ID: 30003339}
\textbf{Source problem code:} catalogue entry
\subsection{Definitions and mathematical statement}
The supplied record concerns historical and contemporary descriptions of successful women mathematicians. It does not define the terms ``gendered descriptions,'' ``successful,'' ``persistence,'' or ``leaky-pipeline phenomenon,'' nor does it specify a population, historical interval, corpus, career stages, measurements, or meaning of correlation.

The retained research question is: In the historical and contemporary context of women in mathematics, how can the persistence of gendered descriptions of successful women mathematicians be analyzed, and does it correlate with the leaky-pipeline phenomenon?

This is a historical, biographical, or sociological research prompt. The supplied evidence does not determine a precise mathematical proposition, statistical hypothesis, domain of quantification, or boundary conditions. Giving such a formulation would require additional assumptions and operational definitions absent from the record. The supplied statement-verification field classifies the mathematical statement as unrecoverable and promotes no separate clean statement; the original question is retained for provenance.
\subsection{Short English statement}
The question asks how persistent gendered descriptions of successful women mathematicians can be studied and whether they are associated with the leaky-pipeline phenomenon. The record supplies insufficient definitions or evidence to formulate a precise mathematical claim. Its supplied ``open'' label describes a research prompt, not an independently verified open mathematical problem.
\subsection{Sources}
\begin{itemize}
\item Local UnsolvedMath registry, permanent ID \texttt{30003339}, title: Gendered Descriptions and the Leaky Pipeline in Mathematics; source code \texttt{catalogue entry}; collection: Oberwolfach Reports Open Problems. The original statement remains retained for provenance. No individual proposer is supplied.
\item Supplied source citation: Mini-Workshop: Women in Mathematics: Historical and Modern Perspectives (2018), Oberwolfach Reports, DOI 10.4171/owr/2017/2, pp.\ 83--131. \url{https://doi.org/10.4171/owr/2017/2}. The record describes Oberwolfach Reports as documenting mathematical workshops and their open-problem discussions. The proposed year and year tag are 2017, whereas the supplied citation gives 2018; this discrepancy is unresolved.
\item Supplied literature source, labeled ``original Women in Mathematics report'': \url{https://doi.org/10.4171/owr/2017/2}.
\item Literature assessment dated exactly 2026-08-21: the supplied curation labels the question open, characterizes it as historical and sociological research about mechanisms behind the leaky pipeline rather than a precise proposition with a settled truth value, and reports no completed answer in the extracted context. This is supplied dated curation derived from the OWR record, not independently verified source consultation or a new literature assessment. Time-sensitive status and cited claims require independent verification.
\end{itemize}
\end{document}
